\documentclass[12pt, a4paper]{article}

\usepackage[utf8]{inputenc}
\usepackage[T1]{fontenc}
\usepackage{geometry}
\usepackage{setspace}
\usepackage{titlesec}
\usepackage{hyperref}
\usepackage{graphicx}
\usepackage{enumitem}

\geometry{margin=1in}
\doublespacing

\hypersetup{
    colorlinks=true,
    linkcolor=blue,
    citecolor=blue,
    urlcolor=blue
}

\title{%
    \vspace{2cm}
    \textbf{Proposal Seed} \\
    \vspace{0.5cm}
    \large CISC-727: Advanced Research Explorations (ARE) - I \\
    \vspace{0.3cm}
    \large Assignment A03
}
\author{%
    Kenneth Peter Fernandes \\
    \vspace{0.3cm}
    Harrisburg University of Science and Technology
}
\date{Spring 2026}

\begin{document}

% ============================================================
% Title Page
% ============================================================
\begin{titlepage}
    \centering
    \vspace*{2cm}
    {\LARGE \textbf{Proposal Seed} \par}
    \vspace{1cm}
    {\large CISC-727: Advanced Research Explorations (ARE) - I \par}
    \vspace{0.5cm}
    {\large Assignment A03 \par}
    \vspace{1.5cm}
    {\large \textit{FlashAttention Speedup Under Multi-GPU Communication Overhead: DDP vs FSDP} \par}
    \vspace{1.5cm}
    {\Large Kenneth Peter Fernandes \par}
    \vspace{0.5cm}
    {\large Harrisburg University of Science and Technology \par}
    \vspace{1cm}
    {\large Spring 2026 \par}
\end{titlepage}

% ============================================================
% Problem Statement
% ============================================================
\section{Problem Statement}

Transformer-based models rely on self-attention, which scales quadratically with sequence length in both time and memory. FlashAttention addresses this bottleneck not by approximating attention but by optimizing memory access patterns at the GPU kernel level---reducing HBM traffic through tiling and fused operations while preserving exact attention outputs. On a single GPU, FlashAttention demonstrates significant speedups over standard attention, particularly at longer sequence lengths.

However, modern training workflows increasingly require multi-GPU parallelism. When training is distributed across GPUs, a new bottleneck emerges: inter-GPU communication (gradient synchronization, parameter sharing). The literature extensively benchmarks FlashAttention in single-GPU and large-cluster settings, but there is a gap in understanding how FlashAttention's compute-side IO-aware advantage interacts with communication overhead in multi-GPU configurations---specifically whether the relative speedup is preserved, diminished, or eliminated as communication burden increases across parallelization strategies like DDP and FSDP.

This gap matters practically: researchers and practitioners with modest multi-GPU setups (2--8 GPUs) need guidance on whether adopting FlashAttention yields meaningful end-to-end training speedups or whether communication overhead dominates total step time, making the kernel-level optimization less impactful than single-GPU benchmarks suggest.

% ============================================================
% Novelty / Positioning Paragraph
% ============================================================
\section{Novelty / Positioning Paragraph}

Existing work on FlashAttention (Dao et al., 2022; Dao, 2023; Shah et al., 2024) demonstrates substantial speedups through IO-aware exact attention on single GPUs, with each successive version improving GPU utilization through better work partitioning and hardware-specific features. Separately, the efficient attention literature (BigBird, Longformer, Performer, Linformer, Nystr\"{o}mformer) pursues sub-quadratic complexity through approximation or sparsity, but cross-family comparisons remain inconsistent due to varying hardware, metrics, and evaluation protocols (Tay et al., 2020; Zhang et al., 2025). What is missing from both lines of work is a controlled investigation of how FlashAttention's compute-side benefit scales under increasing communication overhead in multi-GPU training. This study addresses that gap by measuring FlashAttention's relative speedup across three communication regimes---no communication (single GPU), low communication (DDP), and high communication (FSDP)---on identical hardware, model, and dataset configurations. The contribution is not a new attention mechanism but a reproducible empirical framework that quantifies the interaction between kernel-level compute optimization and distributed training communication cost, providing evidence-based guidance for multi-GPU attention kernel selection.

% ============================================================
% Method Summary
% ============================================================
\section{Method Summary}

The proposed method is an empirical comparative study, not an algorithmic contribution. The approach consists of:

\begin{enumerate}
    \item \textbf{Model and Task:} ViT-Small (patch size 16, embed dim 384, 6 heads, 12 layers) trained on CIFAR-10 for image classification. The model is chosen for its attention-heavy architecture (directly exercises FlashAttention), and the dataset for its simplicity (the research question is about system performance, not model accuracy).

    \item \textbf{Attention Backends:} Each experiment is run with two attention implementations---standard PyTorch math-based attention and FlashAttention via \texttt{scaled\_dot\_product\_attention}---holding all other variables constant.

    \item \textbf{Parallelization Strategies:} Three communication regimes are compared:
    \begin{itemize}
        \item Single GPU (zero communication overhead)
        \item DDP (low communication---all-reduce at end of backward pass)
        \item FSDP (high communication---all-gather per layer + reduce-scatter)
    \end{itemize}

    \item \textbf{Optimization Ablations:} Mixed Precision (AMP) and Overlap Compute/Comm are tested individually and in combination on the DDP + FlashAttention configuration to assess whether compute-side and comm-side optimizations stack with FlashAttention.

    \item \textbf{Measurement:} Time per step, throughput, speedup, scaling efficiency, and compute vs communication breakdown are collected using CUDA event timing. FlashAttention's relative speedup (\%) is computed within each parallelization strategy to test the hypothesis.
\end{enumerate}

% ============================================================
% Evaluation Plan
% ============================================================
\section{Evaluation Plan}

\subsection{Metrics}
\begin{itemize}
    \item Time per step (seconds)---primary timing metric
    \item Throughput (samples/second)
    \item Speedup relative to single-GPU standard attention baseline
    \item Scaling efficiency (speedup / number of GPUs)
    \item FlashAttention relative speedup (\%) within each strategy---primary hypothesis metric
    \item Compute vs communication time breakdown
\end{itemize}

\subsection{Baselines}
\begin{itemize}
    \item Single GPU + Standard Attention (absolute baseline)
    \item Single GPU + FlashAttention (maximum FA benefit reference)
    \item DDP + Standard Attention (multi-GPU baseline without FA)
\end{itemize}

\subsection{Ablations}
\begin{itemize}
    \item 9 runs total: 6 core hypothesis runs (attention $\times$ parallelization) + 3 optimization ablation runs (AMP, overlap, combined)
    \item One-factor-at-a-time discipline: each run changes exactly one variable
\end{itemize}

\subsection{Workloads / Datasets}
\begin{itemize}
    \item CIFAR-10 (pilot). For ARE-II, extend to ImageNet-1K or a long-sequence NLP dataset to test at longer sequence lengths where FlashAttention's advantage is expected to be larger.
\end{itemize}

\subsection{What Constitutes Success}
\begin{itemize}
    \item If FlashAttention's relative speedup (\%) decreases monotonically from single-GPU $\rightarrow$ DDP $\rightarrow$ FSDP, the hypothesis is supported and the communication-overhead modulation effect is confirmed.
    \item If the speedup remains constant, this refutes the hypothesis but provides equally valuable evidence that FlashAttention's benefit is robust to communication regime.
    \item Either outcome produces a publishable finding.
\end{itemize}

% ============================================================
% Pilot Findings (A03)
% ============================================================
\section{Pilot Findings (A03)}

The pilot was executed on Kaggle with 2$\times$ NVIDIA T4 GPUs (PCIe) using ViT-Small (patch 16) on CIFAR-10, running all 9 experiments (R1--R9) with 3 repetitions each.

\subsection{Key Results}

\begin{enumerate}
    \item \textbf{FlashAttention showed no meaningful speedup in any configuration.} Relative speedup was $-2.63\%$ (single GPU), $+1.01\%$ (DDP), and $-0.06\%$ (FSDP)---all within the $\pm 5\%$ noise threshold. FlashAttention and standard attention were functionally equivalent.

    \item \textbf{Multi-GPU training was slower than single GPU.} DDP and FSDP achieved only 0.59$\times$ speedup (vs.\ expected 1.5--1.8$\times$). Communication overhead on PCIe ($\sim$24ms) dominated the $\sim$35ms compute time.

    \item \textbf{AMP increased step time by 9.4\%.} The T4's FP16 overhead exceeded the compute savings for this small model. Overlap compute/comm had negligible effect because there was insufficient compute time to hide communication behind.
\end{enumerate}

\subsection{Interpretation}

The hypothesis (FlashAttention's speedup diminishes with increasing communication overhead) could not be tested because FlashAttention showed no baseline speedup to begin with. The root cause is sequence length: ViT-Small with patch size 16 on 32$\times$32 CIFAR-10 images produces only 5 tokens (4 patches + 1 CLS token). The resulting 5$\times$5 attention matrix fits entirely in GPU SRAM without requiring FlashAttention's tiling---there is no HBM traffic to optimize. This is consistent with the ``efficiency length'' concept from the CAB benchmark identified in the A02 literature review.

This is a valid negative result that directly informs the ARE-II experiment design: the sequence length must be long enough for FlashAttention to demonstrate a baseline speedup before the communication-modulation hypothesis can be tested.

% ============================================================
% Next Steps for ARE-II
% ============================================================
\section{Next Steps for ARE-II}

The pilot identified sequence length as the critical variable that must be addressed before the core hypothesis can be tested. The ARE-II plan is refined accordingly:

\begin{enumerate}
    \item \textbf{Increase sequence length (critical).} Reduce ViT patch size from 16 to 4, increasing sequence length from 5 to 65 tokens on CIFAR-10. Alternatively, switch to an NLP workload with sequence lengths of 512--4096 tokens (e.g., WikiText-103 with GPT-2), where FlashAttention's advantage is well-established in single-GPU benchmarks.

    \item \textbf{Scale hardware.} Move from Kaggle 2$\times$ T4 (PCIe) to a multi-GPU node with NVLink (e.g., 4$\times$ or 8$\times$ A100). The pilot showed PCIe communication overhead dominated step time---NVLink would reduce this bottleneck and allow cleaner measurement of the compute-side FlashAttention effect.

    \item \textbf{Add tensor parallelism.} Include tensor parallelism, which splits the attention computation itself across GPUs, to test whether FlashAttention's tiling interacts with intra-layer parallelism differently than with data-parallel communication.

    \item \textbf{Cross-family comparison.} Extend beyond FlashAttention vs standard attention to include sparse (BigBird) and linear (Performer) alternatives under the same multi-GPU conditions. This addresses the broader A02 gap of inconsistent cross-family evaluation.

    \item \textbf{Formalize the evaluation framework.} Package the measurement methodology (CUDA event timing, compute/comm breakdown, relative speedup metric) as a reusable benchmark protocol that other researchers can apply to new attention kernels and parallelization strategies.
\end{enumerate}

\subsection{Contributions (for ARE-II)}
\begin{itemize}
    \item An empirical finding on how compute-side kernel optimization interacts with distributed training communication overhead
    \item A reproducible evaluation framework for multi-GPU attention kernel benchmarking
    \item Evidence-based guidance for practitioners choosing attention implementations in multi-GPU settings
    \item Characterization of FlashAttention's sequence-length crossover threshold under distributed training conditions
\end{itemize}

\end{document}
